
\begin{theorem}
    Пусть подпространство $\mathcal{E}_1$ инвариантно относительно ортогонального преобразования $\phi$.  Тогда $\mathcal{E}_1^{\bot}$ тоже инвариантно относительно $\phi$.
\end{theorem}
\begin{proof}
    Так как $\phi$ биективно, то $\dim \phi(\mathcal{E}_1)=\dim \mathcal{E}_1$. Но $\phi(\mathcal{E}_1)\subset \mathcal{E}_1$, следовательно $\phi(\mathcal{E}_1)\equiv \mathcal{E}_1$. Осталось заметить, что 
    \[
    \forall x \in \mathcal{E}_1, y\in \mathcal{E}_1^{\bot} \hookrightarrow 0=(x,y)=(\phi(x),\phi(y)) \Rightarrow \phi(y)\in(\phi(\mathcal{E}_1))^{\bot}\equiv\mathcal{E}_1
    \]
    Значит $\phi(\mathcal{E}_1^{\bot})\subset \mathcal{E}_1^{\bot}$.
\end{proof}

\begin{theorem}[о полярном разложении]
    Любое линейное преобразование $A$ евклидова пространства может быть представлено в виде 
    \[
    A=QS
    \]
    где $Q$ -- ортогональное, а $S$ -- самосопряженное с неотрицательными собственными числами.
\end{theorem}
\begin{proof}
    Рассмотрим преобразование $A^*A$: оно самосопряженное (т.к. $(A^*Ax,y)=(Ax,Ay)=(x,A^*Ay)$). Тогда существует (по теореме 13.5) ОНБ из собственных векторов $e^{(n)}$ (соответствующие собственные значения $\lambda_1\geqslant \dots \geqslant \lambda_n$). Тогда
    \[
    (A(e_i), A(e_j))=(A^*A(e_i),e_j)=\lambda_i(e_i, e_j)=\left\{
    \begin{aligned}
    &0 \quad &(i\neq j)\\
    &\lambda_i \quad &(i=j)
    \end{aligned}
    \right. 
    \]
    Значит $\lambda_i=(A(e_i),A(e_j))=|A(e_i)^2|\geqslant 0$ и $|A(e_i)|=\sqrt{\lambda_i}$ -- называется \textit{сингулярным числом} преобразования $A$. 

    Теперь посмотрим на набор векторов $f_i=\frac{A(e_i)}{\sqrt{\lambda_i}}, \ i\in \{1, \dots, k\}$, где $k$ -- последний индекс при котором $\lambda_k \neq 0$. Эти векторы попарно ортогональны и отнормированы.  Дополним их векторами $f_{k+1},\dots, f_n$ до ОНБ $f^{(n)}$. Тогда существует ортогональная матрица перехода $Q$ из $e^{(n)}$ в $f^{(n)}$: $Q(e_i)=f_i$. Определим $S:=Q^{-1}A$. Покажем, что такое преобразование $S$ подходит
    \[
    S(e_i)=Q^{-1}A(e_i)=Q^{-1}\sqrt{\lambda_i}f_i=\sqrt{\lambda_i}e_i
    \]
    Значит собственные значения $S$ неотрицательны, а самосопряженным оно является по теореме 13.5 (т.к. $e^{(n)}$ -- ОНБ из собственных векторов).
\end{proof}

Теперь перейдем в ОНБ ($P$ -- ортогональная матрица перехода) в котором $S$ имеет диагональный вид (т.е. в базис из собственных векторов):
\[
S\to D=P^{-1}SP \Rightarrow S=PDP^{-1}
\]
Тогда преобразование $A$ можно представить в виде
\[
A=QS=QPDP^{-1}
\]
Делая замену $Q'=QP, Q''=P^{-1}$ (по-прежнему ортогональные), получаем \textit{сингулярное разложение}
\[
A=Q'DQ''
\]
Сформулируем это в теорему.
\begin{theorem}[о сингулярном разложении]
    Любое линейное преобразование $A$ евклидова пространства может быть представлено в виде 
    \[
    A=Q'DQ''
    \]
    где $Q',Q''$ -- ортогональные, а $D$ -- диагональное с неотрицательными собственными числами (сингулярными числами).
\end{theorem}

\section{Присоединенные преобразования}
Пусть в евклидовом пространстве $\mathcal{E}$ задана билинейная функция $b(x,y)$.
\begin{definition}
    Линейное преобразование $\phi$ называется присоединенным к функции $b(x,y)$, если 
    \[
    \forall x,y\in\mathcal{E}\hookrightarrow b(x,y)=(x,\phi(y))
    \]
\end{definition}

Докажем корректность данного определения.
\begin{theorem}
    Для любой билинейной формы существует единственное присоединенное преобразование.
\end{theorem}
\begin{proof}
    Пусть $B$ -- матрица билинейной формы $b(x,y)$, а $A$ -- матрица преобразования $\phi$ в некотором базисе. Если $\phi$ является присоединенным, то по определению $x^TBy=x^T\text{Г}Ay$. Отсюда $B=\text{Г}A$, но $\text{Г}$ невырожденная (по теореме 10.3), а значит $A=\text{Г}^{-1}B$ ($\phi$ определяется однозначно). Для доказательства существования заметим, что для любой билинейной формы преобразование с такой матрицей будет присоединенным. 
\end{proof}

\textit{Следствие}: в ОНБ $\text{Г}\equiv E$, а значит и $A=B$.
\begin{theorem}
    Для симметричных билинейных форм и только для них присоединенное преобразование является самосопряженным.
\end{theorem}
\begin{proof}
    Утверждение следует из того, что в ОНБ матрица самосопряженного преобразования симметрична, а матрица присоединенного преобразования равна матрице билинейной формы.
\end{proof}

\begin{definition}
    Преобразование называется присоединенным к квадратичной форме, если оно присоединено к соответствующей билинейной форме.
\end{definition}

\begin{theorem}
    Пусть квадратичная форма $k(x)$ задана в $\mathcal{E}$. Тогда существует ОНБ, в котором $k(x)$ имеет диагональный вид.
\end{theorem}
\begin{proof}
    Таким базисом будет является ОНБ из собственных векторов самосопряженного присоединенного преобразования: в нем матрица присоединенного преобразования имеет вид $diag(\lambda_1, \dots, \lambda_n)$. А по следствию теоремы 15.1 в ОНБ она совпадает с матрицей $k(x)$.
\end{proof}

\subsection{Приведение двух квадратичных форм к диагональному виду}
Постановка задачи: две квадратичные формы требуется привести одновременно к диагональному виду.

\begin{theorem}
    Пусть в линейном пространстве заданы две квадратичные формы $k$ и $h$, причем $h$ положительно определенная. Тогда существует базис в котором обе формы имеют диагональный вид.
\end{theorem}
\begin{proof}
    Пусть изначально $K$ и $H$ -- матрицы $k$ и $h$ соответственно. Определим скалярное произведение формой $h$. Приведем $H$ к каноническому виду, тогда 
    \[
    H\to H'=E\quad K\to K'
    \]
    Из $H'=E$ следует, что полученный базис -- ОНБ (так как $h$ задает скалярное произведение). По теореме 15.3 существует ОНБ из собственных векторов присоединенного преобразования, в котором $K'\to K''=diag(\lambda_1,\dots,\lambda_n)$. Но при переходе к этому базису $H'\to H''$, где $H''$ все еще единичная (так как переход между ОНБ осуществляется с помощью ортогональных матриц: $H''=S^TH'S=S^TES=S^TS=E$). Значит в данном базисе обе формы имеют диагональный вид.
\end{proof}
Более общее утверждение:
\begin{theorem}
    Пусть заданы две квадратичные формы $h$ и $g$. Если среди их линейных комбинаций есть положительно определенная, то существует базис в котором обе формы имеют диагональный вид.
\end{theorem}
\begin{proof}
    Пусть $f=\alpha h+\beta g$ -- положительно определена. Не умаляя общности положим $\beta\neq0$. Тогда по теореме 15.4 существует базис, в котором $f,h$ имеют диагональный вид. Тогда в этом же базисе $g=\frac{f-\alpha h}{\beta}$ тоже имеет диагональный вид.
\end{proof}

\textit{Замечание}: данные условия не являются необходимыми, контрпример: $h(x)=x_1^2-x_2^2$, $g(x)=2x_1^2-2x_2^2$ (обе формы уже имеют диагональный вид, но любая их линейная комбинация не является положительно определенной).

\textit{Замечание}: вообще говоря не любая пара квадратичных форм может быть одновременно приведена к диагональному виду, контрпример: $h(x)=x_1^2-x_2^2$, $g(x)=x_1x_2$.

